\section{Computational methods}
\label{sec:methods}

\subsection{Configuration generation and selection}
\label{sec:configurations}

The reference database was constructed in two stages: exploration of Ti--Al configurations with a pretrained potential, followed by DFT labeling of a compact representative subset.
The medium MACE-MPA-0 foundation model was used for the exploration stage,\cite{Batatia2025} with molecular dynamics and structural optimization implemented through the Atomic Simulation Environment (ASE).\cite{Larsen2017}
The configuration pool combined liquid trajectories, configurations sampled during relaxation from liquid snapshots, and trajectories initiated from substitutional solid solutions.
The crystalline parent structures included hcp and bcc Ti, $\gamma$-TiAl, and $\alpha_2$-Ti$_3$Al.
Off-stoichiometric configurations were generated by randomly replacing atomic species on the parent lattices.
These structures were random substitutional realizations; no optimization of special quasirandom structures was performed.

The solid-solution sampling runs were carried out at 1000~K for 10~ps, whereas the liquid-generation scripts used 2000~K and included alternating MD and relaxation stages.
The exploration dynamics used a 1~fs time step and a target pressure of 1~bar.
Relaxation trajectories supplemented the finite-temperature configurations with lower-energy local environments.
MACE energies and forces were used to generate this pool; the targets for fine-tuning were subsequently obtained from DFT.
% The initially described 1 ps interval between liquid relaxation starts is
% not asserted here: archived script variants use 100 or 500 MD steps.

Representative structures were selected with ASAP\cite{Cheng2020} using smooth overlap of atomic positions (SOAP) descriptors.\cite{Bartok2013}
The SOAP expansion used a radial cutoff of 3.2~\AA, a Gaussian width of 0.4~\AA, four radial basis functions, and angular terms up to $l=3$, with Ti--Al cross-species contributions included.
Atomic descriptors were averaged over each periodic configuration.
Farthest-point sampling in this descriptor space reduced the pool to 1000 configurations.
The selected set contained 880 structures with 216 atoms, 104 with 250 atoms, and 16 with 256 atoms, and included both pure elements and alloy compositions.
% FPS is inferred from the archived ASAP command, which leaves the algorithm
% at the documented default. Record the installed ASAP version before release.

\subsection{Electronic-structure reference calculations}
\label{sec:dft}

Single-point DFT calculations were performed at the selected geometries using VASP\cite{Kresse1996} and the projector augmented-wave method.\cite{Blochl1994,Kresse1999}
Three reference descriptions were considered: PBE, PBEsol, and \rscan.
The supplied PAW datasets were Ti\_pv, with ten valence electrons, and Al, with three valence electrons.
All three reference datasets were generated with a common numerical protocol in VASP~6.3.2 using non-spin-polarized calculations, a plane-wave cutoff of 550~eV, and an electronic convergence threshold of $10^{-6}$~eV.
First-order Methfessel--Paxton smearing with a width of 0.1~eV was employed.
The calculations used the \texttt{Accurate} precision setting, automatic real-space projection, and nonspherical contributions to the PAW exchange-correlation terms.
Explicit configuration-dependent reciprocal-space grids contained 8--12 full-grid points, reduced to 4--6 points in the supplied \texttt{KPOINTS} files.
These grids excluded $\Gamma$ and were identical for a given configuration across all three functionals.
The ions and simulation cells were held fixed during labeling.

The resulting datasets supplied total energies, atomic forces, and virials for fine-tuning.
Their sizes and training/validation partitions are listed in Table~\ref{tab:datasets}.
All retained structures could be matched to the common set of 1000 selected configurations; 992 configurations were present in all three labeled datasets.
There was no configuration overlap between the training and validation subsets of an individual functional.
The partitions were constructed separately, however, and were not identical across functionals.
Accordingly, errors on these validation subsets characterize approximation of the respective DFT datasets and are not treated as errors on a shared independent test set.

\begin{table}[tb]
\caption{Numbers of configurations in the fine-tuning datasets. The selected pool contains 1000 configurations.}
\label{tab:datasets}
\begin{ruledtabular}
\begin{tabular}{lrrr}
Reference & Total & Training & Validation \\
\hline
PBE      & 999 & 850 & 149 \\
PBEsol   & 998 & 849 & 149 \\
\rscan   & 993 & 845 & 148 \\
\end{tabular}
\end{ruledtabular}
\end{table}

\subsection{Fine-tuning DPA4-Neo}
\label{sec:finetuning}

The three Ti--Al potentials were initialized from the pretrained DPA4-Neo model,\cite{Li2026DPA4} using the checkpoint \texttt{DPA4-Neo-OMat24-v20260805.pt}.
Training used the PyTorch implementation in DeePMD-kit.
All three fine-tuning runs used the same hyperparameters, with a 6.0~\AA{} cutoff, two interaction blocks, 32 channels, 16 radial functions, and maximum angular order $l_{\max}=3$.
Energy-derived forces preserve the conservative form of the potential during adaptation.

Low-rank adaptation (LoRA)\cite{Hu2021} was enabled with rank 16 and scaling parameter 16.
The training schedule comprised 50 epochs, with a cosine learning-rate decay from $10^{-4}$ to $10^{-6}$.
The warm-up occupied 0.3\% of the schedule and started at 0.2 times the initial learning rate.
Optimization used HybridMuon with a weight decay of 0.01 and gradient-norm clipping at 5.
The loss combined mean-absolute-error terms for energy, forces, and virial, with constant prefactors of 20, 20, and 5, respectively; the force-norm option was enabled.
Automatic batch sizing was set to \texttt{auto:128}.
The model parameters used single precision, with automatic mixed precision enabled for training.
Source-checkpoint and frozen-export identities are recorded by cryptographic hashes for the held-out evaluation.
% Do not interpret "auto:128" as 128 structures per batch.
% Exact loss normalization / force-norm implementation and optimizer version
% must be pinned to the installed DPA4 code before reproducing the training.

For PBE, the checkpoint at 15\,000 optimization steps (15k) was selected for the updated physical-property calculations as a compromise between training/validation accuracy and the harmonic dynamics of hcp Ti.
Its energy and force errors on the complete training and validation sets are lower than at 5k, while it gives the closest integral agreement with the archived PBE DFT phonon density of states among the five saved checkpoints.
This is a joint error-and-property criterion, not early stopping at a minimum of validation loss: later checkpoints further reduce energy error.
For the DFT-reference benchmark reported here, we use the 20\,000-step checkpoint (20k), which gives smaller energy errors and nearly unchanged force errors relative to 15k.
This choice followed inspection of the test results; it is not a blind post-selection evaluation, although no parameters or energy offsets were fitted to the holdout.
The final PBEsol and \rscan{} checkpoints, at 23\,100 and 23\,200 steps, are used throughout.

The PBE physical-property results use 15k for $\alpha$-Ti phonons, mixing enthalpy, density and excess volume at 1873~K, the complete Ti temperature-dependent density series, and Ti coexistence.
The Ti$_3$Al phonons retain the archived 5k calculations. Figure~\ref{fig:density_ti} retains the complete 5k thermal-density series and adds completed 15k recalculations, with checkpoint identities shown explicitly for this comparison.
Elsewhere we use the label PBE without repeating checkpoint indices; this allocation and the detailed provenance are recorded here and in the Supporting Materials.
Learning curves, complete-set errors, and the phonon diagnostics supporting model selection are also provided there.
The archived PBE DFT phonon diagnostic forms part of selection. The main-text lattice-dynamics comparison uses new, same-functional DFT references and matched DFT/DPA4 calculation settings, together with experimental dispersions.

\subsection{Common held-out DFT benchmark}
\label{sec:testset}

A fixed 50-configuration holdout was selected independently of model errors: 25 configurations from a separate 200-structure liquid archive and 25 from trajectories initialized as hcp, bcc, and $\gamma$-TiAl solid solutions (8, 9, and 8 configurations, respectively).
Selection used composition-balanced farthest-point sampling of normalized, configuration-averaged SOAP descriptors within each family, with a 5~\AA{} cutoff, Gaussian width 0.5~\AA, four radial functions, and $l_{\max}=3$.
No selected geometry matched the supplied training or validation geometries under species-preserving periodic matching, atom permutations, and common translations, using a $5\times10^{-5}$~\AA{} tolerance and fixed cell axes.
An additional near-duplicate screen used a 0.1~\AA{} RMSD threshold together with less than 1\% cell deformation.
This establishes a geometrical holdout from the fine-tuning data; independence of the underlying liquid trajectories and absence of overlap with foundation-model pretraining are not established.
The liquid designation records the source pool, rather than a single equilibrium temperature.

The test configurations were labeled without relaxation using VASP~6.6.1, a 550~eV cutoff, $10^{-6}$~eV electronic convergence, and first-order Methfessel--Paxton smearing of width 0.1~eV.
The PAW datasets were unchanged; reciprocal-space grids were generated with kpLib~1.1.1 using a minimum-distance criterion of 35~\AA.
For each geometry, all three functionals used identical PAW files and reciprocal-space grids.
These test calculations retain the principal numerical settings of the labeling protocol, while the VASP and grid-generator versions differ from those of the original datasets.
Only configurations with converged labels for all three functionals enter a comparison, with geometry and file identities checked before analysis.
The three fine-tuned models were evaluated on the same geometries; the unchanged pretrained model was evaluated once and its predictions compared separately with each functional.
No model parameters or energy offsets were fitted to the holdout.

Energy errors compare raw model energies with VASP free energies (\texttt{TOTEN}), divided by the number of atoms.
Force MAE and RMSE pool all Cartesian components; stress errors pool the six independent tensor entries in GPa, with model stress obtained as minus virial divided by volume.
Descriptive 95\% percentile intervals use 2000 bootstrap resamples of whole configurations, retaining the force components within each configuration together and using matched draws across functionals.
These intervals are conditional on the currently completed subset and do not account for unknown trajectory correlations or repeated inspection as the test set grows.

\subsection{Thermal density and liquid mixing thermodynamics}
\label{sec:thermodynamics}

Density and mixing thermodynamics were evaluated using periodic NPT simulations in ASE with the fine-tuned potentials.
The integration time step was 1~fs and the external pressure was 1~bar.
The stored MD implementation used the Melchionna NPT integrator, a thermostat time parameter of 25~fs, and isotropic cell fluctuations.
The barostat parameter was set to $B\tau^2$, where $\tau=25$~fs and $B$ was the bulk-modulus estimate supplied by the workflow.
For liquid Ti--Al mixtures, this estimate was interpolated between 110~GPa for Ti and 76~GPa for Al.
% The barostat interpolation values are checked against utils/bulk_modulus.py.

For the temperature-dependent density of Ti, 250-atom cells were initialized from $\alpha$-Ti at 300, 500, 700, 900, and 1100~K and from $\beta$-Ti at 1200, 1400, 1600, 1800, and 1900~K.
The stored solid-state runs comprised 20~ps of equilibration and 30~ps of production.
The original liquid-Ti series used 200-atom cells between 2000 and 2500~K in 100~K increments.
Its liquid preparation involved 10~ps at $1.5T$, followed by 30~ps at the target temperature and a further 30~ps of production.
Production configurations were saved every 0.1~ps.
The initial crystalline phase labels specify the preparation procedure; interpretation of a high-temperature density point also requires checking whether that phase persisted during the trajectory.

For the additional PBE 15k temperature scan, the liquid starting cell is taken from the completed 2000~K production trajectory and equilibrated independently for 30~ps at each target temperature, followed by 30~ps of production.
The 250-atom $\alpha$ and $\beta$ starting cells are built from their respective 15k static equilibrium lattice parameters and use 20~ps of equilibration followed by 30~ps of production.
Velocities and thermostat variables are reset at the production boundary. Scalar observables are saved every 10~fs and configurations every 0.1~ps.
Only completed production windows enter Fig.~\ref{fig:density_ti}; density is the mean instantaneous mass divided by volume, and its SEM for this PBE 15k scan is the largest estimate from 3, 5, and 6~ps blocks.
The thermostat targets $3(N-1)$ kinetic degrees of freedom, whereas the unconstrained ASE temperature uses $3N$; the logged temperature is corrected for this distinction when checking the target.

The PBEsol and \rscan{} liquid-Ti branches in Fig.~\ref{fig:density_ti}(b) use extended sampling at 1873~K and at 2000--2500~K in 100~K increments.
Each 200-atom cell starts from the final archived liquid configuration of the same functional and temperature, followed by 10~ps re-equilibration and 90~ps new production with the same NPT settings and $B=110$~GPa.
Velocities and thermostat/barostat variables are reset at each stage; these runs are new production windows, not exact trajectory continuations.
The mean uses 9000 scalar samples at 10~fs spacing, excluding the initial frame; configurations are saved every 100~fs. The archived 30~ps data are not pooled with the new estimates.
The SEM is the maximum over 3, 5, 6, 9, 10, and 15~ps blocks, requiring at least six blocks. No fit or choice of averaging window is used to smooth the curves.
Phase checks combine non-affine, center-of-mass-corrected 10~ps displacements with local-order diagnostics every 2~ps, including the final frame; details are given in the Supporting Materials.
The mean pressure is compared with the 1~bar target using a block Student-$t$ 95\% interval. Residual pressure offsets are reported as a protocol limitation and are not included in the density SEM.

For the alloy liquids, we define the composition as the Al atomic fraction,
\begin{equation}
 x = x_{\mathrm{Al}} = \frac{N_{\mathrm{Al}}}{N_{\mathrm{Al}}+N_{\mathrm{Ti}}}.
\end{equation}
Alloy-liquid calculations used 200 atoms at $x=0,0.1,\ldots,1$.
The composition-dependent density and excess-volume comparison uses 1873~K, whereas the mixing-enthalpy comparison uses 2000~K.
Density and mixing enthalpy were originally obtained from the same production trajectories.
For the PBE mixing-enthalpy series at 2000~K, each of the 11 compositions was initialized from the final configuration of a previously equilibrated liquid and propagated for 30~ps of equilibration plus 30~ps of production.
Velocities and the NPT integrator were reinitialized at the stage boundary.
Scalar observables were recorded every 10~fs and full configurations every 100~fs.
The 30~ps production interval is the primary estimate; an extended 49~ps window is assessed only as a sensitivity check in the Supporting Materials.
The PBE density and excess-volume series at 1873~K uses the same 30~ps equilibration plus 30~ps production protocol for all 11 compositions, initialized from the corresponding archived 1873~K liquids.
The pure-Ti endpoints at 1873 and 2000~K also enter the temperature-dependent liquid-density figure.
These two updated density estimates use the 10~fs scalar logs rather than the 100~fs configuration sampling used for the retained series.
The mass density was calculated as the time average $\rho=\langle M/V\rangle$, where $M$ is the total mass of the simulated cell.

For each model, the molar enthalpy was evaluated as
\begin{equation}
 H_m(T,p,x)=\frac{N_{\mathrm A}}{N}
 \left[\langle E_{\mathrm{pot}}+E_{\mathrm{kin}}\rangle+
 p_{\mathrm{ext}}\langle V\rangle\right],
 \label{eq:enthalpy}
\end{equation}
where $N_{\mathrm A}$ is Avogadro's constant and $N$ is the number of atoms.
The mixing enthalpy is
\begin{align}
 \Delta H_{\mathrm{mix}}(T,p,x)={}&H_m(T,p,x)\nonumber\\
 &-xH_m^{\mathrm{Al}}(T,p)-(1-x)H_m^{\mathrm{Ti}}(T,p).
 \label{eq:mixing}
\end{align}
The pure-liquid reference values were evaluated with the same potential at the same temperature and pressure.
Equation~\eqref{eq:enthalpy} uses the prescribed external pressure; the ASE stress convention is $p=-\mathrm{Tr}\,\boldsymbol{\sigma}/3$.
The enthalpies were reprocessed from the stored energies and volumes using this definition.
At a common temperature the classical kinetic contribution cancels in the thermodynamic mixing limit.
The excess molar volume was calculated analogously from $V_m=N_{\mathrm A}\langle V\rangle/N$:
\begin{equation}
 V_m^{\mathrm{ex}}(x)=V_m(x)-xV_m^{\mathrm{Al}}-(1-x)V_m^{\mathrm{Ti}}.
\end{equation}
Liquid-state values below the corresponding liquidus describe undercooled liquids and must be distinguished from stable equilibrium phase mixtures.

Statistical variability was assessed by subdividing production trajectories into contiguous time blocks.
The PBEsol enthalpy estimates and the retained density series use six 5~ps blocks.
For \rscan{} mixing enthalpy at 2000~K, every composition and both pure-component references were extended by 30~ps from the final coordinates, cell, and momenta, giving 60~ps of accumulated production.
The duplicate boundary frame was excluded; thermostat and barostat variables were reinitialized at the continuation.
The main estimate uses twelve 5~ps blocks, with block-length and restart sensitivity assessed in the Supporting Materials.
For the PBE series, the conservative standard error of each trajectory mean is the maximum obtained with 3, 5, and 6~ps blocks, requiring at least four blocks; these errors are propagated to the mixing enthalpy and excess molar volume.
Additional block lengths and averaging windows are examined in the Supporting Materials.
Uncertainty in a mixing quantity includes the uncertainties of both pure-component reference values; reuse of these references also correlates errors between compositions.
The final uncertainty estimates require stabilization with respect to block length and sampling time.

\subsection{Harmonic lattice dynamics}
\label{sec:phonons}

For $\alpha$-Ti, direct DFT and fine-tuned DPA4 phonons were evaluated with a common finite-displacement protocol in Phonopy.\cite{Togo2015}
The hcp unit cell was relaxed separately with PBE, PBEsol, and \rscan{}; within each functional, DFT and DPA4 forces were then evaluated on exactly the same displaced structures, without a separate DPA4 relaxation.
The force-constant supercell was $4\times4\times3$ (96 atoms), with symmetry-generated opposite displacements of $\pm0.01$~\AA{} and a Ti mass of 47.867~amu.
The DFT force calculations used VASP~6.6.1, the Ti\_pv PAW dataset, a 550~eV cutoff, a $\Gamma$-centered $6\times6\times6$ electronic grid, Gaussian smearing of width 0.1~eV, and an electronic convergence criterion of $10^{-8}$~eV.
Real-space projection was disabled and nonspherical PAW contributions were included.
For both methods, mean force drift was removed and the central difference of the opposite-displacement forces was used to construct force constants, followed by identical symmetry projection and enforcement of the acoustic sum rule.
Thus, differences within a DFT/DPA4 pair probe the model's harmonic response at fixed geometry and numerical phonon settings; differences between the three DFT spectra also include the functional dependence of the equilibrium cell.

Dispersions were evaluated along $\Gamma$--K--M--$\Gamma$--A and compared with 77 digitized room-temperature experimental points reproduced in Ref.~\onlinecite{Souvatzis2007}.
The density of states (DOS) used a $\Gamma$-centered $50^3$ phonon mesh, Gaussian broadening $\sigma=0.05$~THz, and a 0.002~THz frequency spacing; each DOS integrates to three modes per atom.
The calculations are harmonic spectra at static relaxed geometries and exclude thermal expansion and anharmonic frequency renormalization.
DFT cutoff, electronic-grid, smearing, and supercell convergence are not yet established; the current spectra and their comparisons are conditional on the stated settings.
Reference-force checks and the cancellation of common reference forces in the central difference are documented in the Supporting Materials, together with the earlier diagnostics used for model selection.

The retained $\alpha_2$-Ti$_3$Al model calculations use ASE with 1000-atom supercells.
Their BFGS relaxation preserves crystal symmetry and allows hydrostatic cell strain, with a force convergence target of $10^{-4}$~eV\,\AA$^{-1}$.
Positive and negative Cartesian displacements of 0.01~\AA{} were used to construct force constants, followed by enforcement of the acoustic sum rule.
No corresponding DFT phonon benchmark is included for Ti$_3$Al.

\subsection{Solid--liquid coexistence and melting of Ti}
\label{sec:melting}

The melting temperature of Ti was estimated by direct solid--liquid coexistence simulations in LAMMPS.\cite{Thompson2022}
A periodic 1296-atom cell, approximately $20\times20\times59$~\AA$^3$ at the beginning of production, contained a bcc solid slab and a liquid slab separated along $z$.
Periodic boundaries produced two interfaces.
The preparation included 5 or 10~ps of initial equilibration, 20~ps of heating one half of the sample to $1.5T$, and 10~ps of cooling that region to the target temperature while retaining a crystalline seed.
The complete system was then propagated in the NPT ensemble at zero external pressure.
The PBE series uses 10+20+10~ps preparation and 50~ps production at 1700, 1750, 1800, and 1850~K; the additional 1775~K trajectory extends production to 100~ps.
The retained PBEsol and \rscan{} production trajectories are 50~ps long.
The time step was 1~fs; thermostat and barostat damping times were 0.1 and 1~ps, respectively.
The three box lengths were allowed to fluctuate independently.
Configurations were saved every 0.1~ps for PBE and every 0.5~ps for the retained trajectories; structural analysis uses a common 0.5~ps interval.

The direction of interface motion was determined from spatially resolved structural order and the evolution of the solid-like fraction.
For a quantitative check, neighbor-averaged bond-order parameters $\bar q_6$ were calculated following Lechner and Dellago,\cite{Lechner2008} using the 14 nearest neighbors.
A threshold of $\bar q_6>0.25$ was used to identify solid-like atoms, with alternative thresholds and a bond-correlation criterion used to assess classification sensitivity.
Profiles were evaluated along the reduced coordinate $z/L_z$ to separate interface motion from changes in box length.
Growth and recession of the crystalline layer were used to bracket the melting temperature, with complete crystallization or melting providing additional endpoint checks.
The resulting temperature intervals describe the observed behavior of this finite cell over 50~ps; their half-widths are not interpreted as full statistical uncertainties.

\subsection{Liquid viscosity}
\label{sec:viscosity}

Shear viscosity was calculated from equilibrium stress fluctuations using the Green--Kubo relation.
Simulations with the PBE- and \rscan{}-based potentials used 512-atom periodic cells, NPT preparation at the target temperature and 1~bar, and fixed-volume NVT production with a 1~fs time step.
The composition dependence was evaluated at 1873~K for eleven compositions spanning pure Ti to $x_{\mathrm{Al}}=0.90$; pure Al was not included.
The temperature dependence of pure liquid Ti was examined between 1873 and 2100~K.

For an isotropic liquid, define five equivalent shear channels,
\begin{align}
 Q_1&=P_{xy},\quad Q_2=P_{xz},\quad Q_3=P_{yz},\nonumber\\
 Q_4&=\tfrac12(P_{yy}-P_{zz}),\nonumber\\
 Q_5&=\frac{P_{xx}-\tfrac12(P_{yy}+P_{zz})}{\sqrt{3}}.
\end{align}
The viscosity of a fixed-volume trajectory is
\begin{equation}
 \eta=\frac{V}{5k_{\mathrm B}T}\sum_{a=1}^{5}
 \int_0^\infty
 \left\langle\delta Q_a(0)\delta Q_a(t)\right\rangle\,dt,
 \label{eq:viscosity}
\end{equation}
where $\delta Q_a=Q_a-\langle Q_a\rangle$ and the pressure tensor includes both kinetic and configurational contributions.
The integral was estimated with STACIE (STable AutoCorrelation Integral Estimator),\cite{Toraman2025} which obtains viscosity from the low-frequency stress spectrum.
The five channels were combined with the appropriate volume and temperature prefactors for each trajectory.
Reported errors are estimator standard errors; they do not include all effects of finite sampling, preparation, finite cell size, or the potential approximation.

The pure-Ti temperature dependence was summarized by the Arrhenius form $\eta(T)=\eta_0\exp(B/T)$, fitted to the logarithms of the calculated viscosities with weights $(\eta/\sigma_\eta)^2$, where $\sigma_\eta$ is the reported standard error.
These fits describe the sampled temperature range and leave the individual viscosity estimates and their errors unchanged.

For comparison, composition-dependent viscosities were also calculated using the Kozlov--Romanov--Petrov (KRP) relation,\cite{Kozlov1983}
\begin{equation}
 \ln\frac{\eta}{\eta_{\rm ref}}=(1-x)\ln\frac{\eta_{\rm Ti}}{\eta_{\rm ref}}
 +x\ln\frac{\eta_{\rm Al}}{\eta_{\rm ref}}
 -\frac{\Delta H_{\rm mix}}{3RT},
 \label{eq:krp}
\end{equation}
where $x=x_{\rm Al}$, $R$ is the molar gas constant, $\eta_{\rm ref}=1$~mPa\,s, and $\Delta H_{\rm mix}$ is the mixing enthalpy per mole of atoms.
All three KRP curves use $T=1873$~K and the same pure-component viscosities: $\eta_{\rm Ti}=3.457$~mPa\,s from the Ishikawa correlation retained in Fig.~\ref{fig:viscosity}b,\cite{Ishikawa2012} and $\eta_{\rm Al}=0.497$~mPa\,s from the critically evaluated experimental correlation of Assael \textit{et al.},\cite{Assael2006} $\log_{10}[\eta_{\rm Al}/(\mathrm{mPa\,s})]=-0.7324+803.49\,\mathrm{K}/T$.
The Al value is an extrapolation beyond that correlation's 933--1270~K range; its reported approximately 14\% uncertainty is not established at 1873~K.
The enthalpy inputs are, separately, the calorimetric digitization reproduced in Ref.~\onlinecite{Witusiewicz2008} and the archived PBE and \rscan{} MD results at 1873~K, using the same potential parameter sets as the direct viscosity estimates in Fig.~\ref{fig:viscosity}a.
For the calorimetric input only, $\Delta H_{\rm mix}(1873\,\mathrm{K})\simeq\Delta H_{\rm mix}(2000\,\mathrm{K})$ is assumed.
Mixing enthalpies are linearly interpolated between sampled compositions with zero values at the pure-component limits; no parameters are fitted to alloy-viscosity data.
